\documentclass{article}
\usepackage{graphicx} % Required for inserting images
\usepackage{xcolor} % Must be loaded before soul
\usepackage{soul} % For highlighting text
\usepackage[utf8]{inputenc}
\usepackage{newunicodechar}
\usepackage{tikz}
\usetikzlibrary{positioning}
\usepackage[utf8]{inputenc}
\usepackage[T1]{fontenc}
\usepackage{fontspec}           % For Greek text support with XeLaTeX/LuaLaTeX
\usepackage{polyglossia}        % Multilingual support
\setmainlanguage{english}
\setotherlanguage{greek}

\usepackage{geometry}
\geometry{margin=1in}

\usepackage{csquotes}           % Block quotes
\usepackage{enumitem}           % Better lists
\usetikzlibrary{arrows.meta, positioning, shapes.geometric}

\usepackage{fancyvrb}           % Verbatim environments for ASCII diagrams
\usepackage{setspace}           % Line spacing



% --- Greek support (XeLaTeX) ---
\usepackage{fontspec}
\newfontfamily\greekfont{Gentium Plus} % good polytonic Greek; try Times New Roman if needed
\newcommand{\gk}[1]{{\greekfont #1}}

\newunicodechar{₁}{$_1$}

\usepackage[style=mla]{biblatex} % MLA Citation Style for Bibliography
\addbibresource{references.bib} %Add References page for in text citation
\usepackage[colorlinks=true, linkcolor=blue, urlcolor=blue, citecolor=blue]{hyperref} %enable hyperlinks to be colored and interactable
\usepackage{changepage} %create individual blocks of text that can be altered apart from general document.

\newcommand{\ra}{\ensuremath{\rightarrow}}

\newcommand{\inlinenote}[1]{\par\noindent
  \fcolorbox{orange}{orange!20}{\parbox{0.97\linewidth}{#1}}\par\noindent}

\begin{document}

\inlinenote{phantasia presents an object as to-be-pursued or to-be-avoided (practical seeing-as) Nussbaum. Frede's careful analysis of the cognitive role of \textit{phantasia} further supports the dual-trace thesis by noting that the ``uncontrolled status of \textit{phantasiai} as after-images'' does not do justice to their role ``in memory, thinking, or decision-making'' (Frede, \textit{The Cognitive Role of Phantasia in Aristotle}, pp.~4--5).}

\subsection*{A\textsubscript{3}: Completed Cognitive Actuality and the Three Orientational Modes}

The actualization of perception, as I established in the preceding section, terminates not in an ephemeral event, but in a residual composite consisting of dual-traces---a \textit{resonant kinēsis} constituted by  \textit{resonant aisthēma} (formal-epistemic content) and \textit{resonant epithymia} (affective valence)---that persists after the sensible object has withdrawn. This residual substrate, still kinetic and still bearing the formal and affective imprint of the perceived particular, constitutes the $A_1$ moment: perception completed but not erased, a movement that endures without an external mover sustaining it. Now we must address the motion from $M_2 \to A_3$, the passage from that persisting yet indeterminate \textit{resonant kinēsis} to a determinate presentational image available for cognitive operations: the \textit{phantasma} proper, which I designate as occupying $A_2$ in the actualization chain. Aristotle, as we have seen, situates \textit{phantasia} as a capacity that stands between bare perception and rational thought, thereby making the transition from $A_2 \to A_3$ the hinge upon which the entire cognitive architecture turns. Indeed, as I will demonstrate, \textit{phantasia} occupies an architectonic position insofar as it orders the transition from aisthēsis to thought (\textit{nous}), furnishing the \textit{phantasmata} through which \textit{nous} and \textit{doxa} can engage what sensation leaves behind.\footnote{Aristotle establishes the architectonic dependence at \textit{De Anima} III.7, 431a14--17: ``to the thinking soul, images serve as if they were contents of perception. When the mind actively thinks, it is necessarily aware of something along with an image,'' and ``the soul never thinks without an image'' (431a16--17). The claim is repeated at III.8, 432a8 and 432a14, marking it as one of the strongest theses in the \textit{De Anima} on \textit{noēsis}'s structural dependence on \textit{phantasia}.} Heidegger, similarly, reads Aristotle's account of perception as fundamentally a matter of \textit{pathein}---being ``affected by beings as world---wherein the perceiver stands in a relation of being-in-the-world that is already structured by disposition and orientation (\textit{BCAP} 265).\footnote{``this genuine \gk{ἀντί} is: the to-which of being-related in person, i.e., to show oneself, showing in the mode of being-uncovered, there in discovered- (265) ness, i.e., there for ... and to perceive it as one, \gk{πάσχειν}, affected by beings as world; \gk{πάσχειν} as \gk{δέχεσθαι}, ``to perceive'' and that primarily as \gk{πάθος}, ``becoming-affect,'' ``disposition,'' being related to the world as \textit{in} it, being in dealings with it'' (264-265).}\footnote{``Being-there as living, a being as being toward the world. Human being-there determined by \gk{νοῦς}; discoveredness is in this \textit{being-supposing}. This supposing as basic mode of human being-in-the-world is ultimately also a perceiving-it, having-it-there-thus-and-so, i.e., of the world, not only determinate beings in determinate kinds of encounters, but \textit{becoming-affected} by all possible beings of the world (cf. \gk{ἡδονή---καταφάναι, ἀποφάναι---διά}). This \gk{νοεῖν} has the being-character of becoming-affected by what is discovered, what, for its part, is only possible in such a way that this becoming-affected by what is discovered is grounded in a \textit{generally-being-discovered} and a \textit{being-discovered}, i.e., in a discovering, a giving-sight-of \textit{as such}'' (\textit{BCAP} 264-265).} The passage from $A_2 \to A_3$ thus marks the moment at which the soul's being-affected hardens into a determinate presentation, a movement that is at once physiological and cognitive.\footnote{Michalski reads Heidegger's Aristotelian seminars as a hermeneutic-phenomenological philology that ``exceeds the horizon of classical philology'' (Michalski 72): philology, in Heidegger's sense, is ``passion for the knowledge of that which has been expressed and of its self-expression'' (GA 18, p.~4) --- the necessary element of hermeneutic phenomenology. The architectonic position of \textit{phantasia} the present paragraph defends is the philological-philosophical site at which the Aristotelian text is most precisely read through Heidegger's GA 18 idiom: the cognitive architecture is itself a structured manner of being-affected, and the \textit{phantasma}-as-determinate-presentation is the moment at which this being-affected acquires its specifically intelligible form.}

\textit{Phantasia} names the soul's capacity and activity of generating presentation images from perceptual residue; it is the faculty ``that in virtue of which an image arises for us ... in virtue of which we discriminate and are either in error or not (\textit{On the Soul} III.3, 428a1-3). The \textit{phantasma}, on the other hand, names the determinate image that results from that activity. This distinction is well articulated in secondary literature\footnote{Frede argues that \textit{phantasia} is supervenient on but irreducible to sensation, with \textit{phantasmata} functioning as the cognitive vehicle of all post-sensory and non-perceptual thought (Frede 1992, \textit{The Cognitive Role of Phantasia in Aristotle}, pp.~4--5, 279--282). Papachristou refines the picture by distinguishing \textit{three} kinds or grades of \textit{phantasia} --- indeterminate (\textit{aoristos}), sensitive (\textit{aisth\=etik\=e}), and calculative/deliberative (\textit{logistik\=e}/\textit{bouleutik\=e}) --- corresponding to imperfect animals, perfect animals with multiple senses, and rational animals respectively (Papachristou 2013, ``Three Kinds or Grades of \textit{Phantasia} in Aristotle's \textit{De Anima},'' \textit{Journal of Ancient Philosophy} 7.1: 19--48, esp.~pp.~19--27). The two interlocking distinctions --- faculty/product (Frede) and tripartite-grade taxonomy (Papachristou) --- jointly underwrite the \textit{phantasia}/\textit{phantasma} distinction this section deploys throughout.}, and is essential for understanding how Aristotle's cognitive psychology operates. Papachristou insists on leaving the Greek word \textit{phantasia} untranslated precisely because its standard English rendering as ``imagination'' obscures the tripartite ambiguity of the term---its simultaneous reference to capacity, activity, and product (28 footnote 32). Indeed, to translate \textit{phantasia} simply as ``imagination'' would collapse the ontological distinction between the soul's active power of image-formation and the image that stands before the mind as a result of that power. Aristotle's account requires us to recognize that \textit{phantasia} admits of internal grades: specifically, the distinction between sensory \textit{phantasia}, which operates in all animals possessing perception,\footnote{Aristotle distinguishes the two higher kinds of \textit{phantasia} at the close of \textit{De Anima} III.10--11: ``all \textit{phantasia} is either calculative or sensitive'' (III.10, 433b29); the sensitive grade ``exists also in the other animals'' while the calculative or deliberative grade ``exists in those animals which can calculate'' (III.11, 434a5--10), distinguished by the cognitive operation that ``one has the ability to make one \textit{phantasma} out of many \textit{phantasmata}'' (\gk{ἓν ἐκ πλειόνων φαντασμάτων ποιεῖν}, 434a9--10).} and deliberative \textit{phantasia}, which belongs only to rational animals capable of calculation and planning. Sensory \textit{phantasia} generates \textit{phantasmata} that are direct descendants of perceptual episodes, bearing the formal character of the sensible particular from which they derive.\footnote{Aristotle grounds the derivation at \textit{De Anima} III.3, 428b10--17: ``imagination is held to be a movement and to be impossible without sensation, i.e.\ to occur in beings that are percipient and to have for its content what can be perceived, and since movement may be produced by actual sensation and that movement is necessarily similar in character to the sensation itself, this movement cannot exist apart from sensation or in creatures that do not perceive.''} Deliberative \textit{phantasia}, by contrast, synthesizes multiple \textit{phantasmata} into composite images that represent possible future states of affairs. \textit{Phantasia}, to reiterate, as an activity, produces the \textit{phantasma} as product, and the two are related as the act of stamping is related to the impression in wax. Each \textit{phantasma} is a singular image derived from a particular perceptual encounter and participates in the process that make it available for cognition. 

The \textit{phantasma} proper, as constituted at $A_2$, exhibits three essential characteristics that distinguish it from both the \textit{resonant kinēsis} (raw perceptual residue) of $A_1$ and the conceptual content that serves as the foundation for intellection, memory, deliberative and speculative reasoning, and \textit{doxa} formation. First, the \textit{phantasma} is determinate: it possesses a definite formal structure that allows it to present a particular content to the mind, rather than remaining an inchoate sensory afterimage. Second, the \textit{phantasma} is available: it stands before the cognitive faculties as an object that can be contemplated, recalled, combined, or judged,\footnote{The three characteristics --- determinacy, availability, representationality --- track three Aristotelian features articulated at \textit{De Anima} III.3, 427b14--429a9: (i) the \textit{phantasma}'s derivation from sensation as a quasi-causal residue (\textit{kin\=esis} \textit{hypo t\=es energeias t\=es aisth\=eseos}, 428b13); (ii) its capacity for falsehood, where sensation is always true (428a11--12; 428b17--18); (iii) its freedom from belief-conviction, since one can imagine fearful things without being affected (427b21--24). The recent literature converges in reading the \textit{phantasma} as a structurally pre-conceptual vehicle of intentionality: Frede (1992) treats \textit{phantasia} as supervenient on but irreducible to sensation; Caston (1996) reads Aristotle as introducing \textit{phantasia} precisely to solve the \textit{problem of error}, with the \textit{phantasma} functioning as ``a second, causally-mediated change'' (Caston 1996, \textit{Why Aristotle Needs Imagination}, p.~47, esp.~pp.~26--52). The three characteristics are thus not extrinsic features but the structural articulation of \textit{phantasia}'s mode of being as causally-derived-but-irreducible-to-sensation.} thus serving as the common substrate for such cognitive operations. Third, the \textit{phantasma} is representational: it bears a structural likeness to the sensible particular from which it derives, and this likeness is what allows it to function both as an object of contemplation in its own right and as a sign pointing beyond itself to the thing it resembles. These three characteristics---determinacy, availability, and representationality---are not extrinsic features but follow from the essential nature of \textit{phantasia} itself as a movement resulting from actual sensation. Aristotle makes clear that the soul never thinks without a \textit{phantasma},\footnote{The thesis is Aristotle's most-quoted statement on the \textit{phantasma}-grounding of thought: ``the soul never thinks without an image'' (\textit{De Anima} III.7, 431a16--17; \gk{οὐδέποτε νοεῖ ἄνευ φαντάσματος ἡ ψυχή}); cf.~the repetitions at 432a8 and 432a14. Heidegger develops the same claim at \textit{Basic Concepts of Aristotelian Philosophy} p.~134: ``insofar as \textit{noesis} is the highest possibility for the being of human beings, the entire being of human beings is determined so that it must be apprehended as the \textit{bodily being-in-the-world}'' --- i.e., \textit{noein} itself must be grounded in \textit{phantasia}, since the soul thinks always through and with the \textit{phantasma}.} and this dependence of higher order cognitive functions upon the image presupposes that the image possesses sufficient formal determinacy to bear the intelligible structure on which such functions depend. Burke adds a complementary rhetorical perspective in his analysis of the traditional understanding of ``imagination'' when he notes William Hazlitt's use of \textit{identification}\footnote{Burke's canonical definition of \textit{identification} grounds the analysis: ``In pure identification there would be no strife. Likewise, there would be no strife in absolute separateness, since opponents can join battle only through a mediatory ground that makes their communication possible, thus providing the first condition necessary for their interchange of blows. But put identification and division ambiguously together, so that you cannot know for certain just where one ends and the other begins, and you have the characteristic invitation to rhetoric. Here is a major reason why rhetoric, according to Aristotle, `proves opposites''' (\textit{RoM}, p.~25). The identification-division dyad is the structural condition for rhetoric itself; the Hazlitt deployment elaborates this in the specific case of appetite-imagination interaction, where the self's identification with future-self via imagination (\textit{RoM}, pp.~83--84, quoted below) instantiates the more general identification-division mediation that Burke names \textit{par excellence} the topic of a rhetoric ``having `identification' as its key term'' (\textit{RoM}, p.~24).} to account for the interaction of appetite and imagination when the motive of self-interest supersedes the interest of others.\footnote{``The only reason for my preferring my future interest to that of others must arise from my anticipating it with greater warmth of present imagination. It is this greater liveliness and force with which I can enter into my future feelings, that in a manner identifies them with my present being: and this notion of identity being once formed, the mind makes use of it to strengthen its habitual propensity, by giving to personal motives a reality and absolute truth which they can never have. Hence it has been inferred that my real substantial interest in anything must be derived from the impression of the object itself, as if that could have any sort of communication with my present feelings, or excite any interest in my mind, but by means of the imagination, which is naturally affected in a certain manner by the prospect of future good or evil'' (qtd. in \textit{RoM}, 83-84).} This insight, highlighted by Burke, reveals that the \textit{phantasma}'s representational character extends beyond the epistemological to the motivational: the image does not merely present content for contemplation but engages the \textit{orectic} dimension of the soul, moving the desire toward or away from what the image represents. 

Heidegger’s phenomenological reading of De Anima III.3–5 understands \gk{φαντασία} as a \gk{ποιεῖν παρόν}, a ‘making‑present’ that takes up the \gk{εἶδος} or ‘look’ of what has been perceived and holds it in availability for \gk{δόξα}, memory, and deliberation. As such, the phantasma is essentially determinate (it presents a being ‘as such‑and‑such’ and so can be true or false), available (it furnishes the very there in which we remember and think), and representational (it bears the same look as the sensible object and thereby can move us affectively toward or away from it). On this reading, Aristotle’s claim that ‘the soul never thinks without a phantasma’ does not merely add a psychological condition to intellection but indicates that all higher cognition is grounded in this image‑based mode of being‑in‑the‑world.

Heidegger's phenomenological reading of \textit{On the Soul} III.3–5\footnote{See \textit{BCAP} 132-137.} understands \gk{φαντασία} (\textit{phantasia}) as a \gk{ποιεῖν παρόν}, a ``making-present'' that takes up the \gk{εἶδος} (\textit{eidos}) or ``look'' of what has been perceived and holds it in availability for \gk{δόξα} (\textit{doxa}), memory, and deliberation. The decisive move, however, is that this making-present is not the activity of a faculty producing inner copies of external objects; it is a \textit{mode of being-in-the-world}, ``the ``making-present-to-itself'' of the world, in which what is made present is not actually there, but instead is, say, in memory or in a merely faint making-present'' (134). The formulation rejects, in advance, the picture of \textit{phantasia} as image-formation in a private mental theater, and it correspondingly transforms what would otherwise look like three features of an inner mental image---determinacy, availability, representationality---into three structural moments of the soul's disclosedness toward the \textit{eidos} of a being.

Determinacy is, on this reading, the \textit{eidos}-character of the look itself: \textit{phantasia}, through the \textit{phantasma}, presents a being as-such-and-such, and so can be true or false, because the look it has is the look of \textit{this} being and not another. Availability is the very 'making-present' that grounds remembering, thinking, and deliberation---the furnishing of the 'there' in which a matter can be held before us when it is not perceptually present. And representationality is the \textit{uncovering} function of \textit{phantasia}: in bearing the same look as the sensible object, the \textit{phantasma} makes the thing itself manifest in a different, non-sensory mode---not merely epistemically, but as the mode in which a possible good or bad is made present \textit{as mattering for me}. What \textit{phantasia} makes present, in each of these moments, is the \textit{eidos}---the look of a being, the ``what it is''---and the hedonic tonality or \textit{resonant epithymia}---the ``how it matters to me.'' \textit{Phantasia} is, therefore, the light in which the look of something is seen and felt: ``Just as it is through light that a color first comes to its being-there, is in its there insofar as it stands in illumination---being-there as the characteristic illumination, so \textit{every being that is there as a being requires a fundamental illumination in order to be there}.... This possibility is nothing other than \gk{νοῦς} [\textit{nous}]'' (135). The \textit{phantasma} is thus the way in which the \textit{eidos} and hedonic tonality of a being is taken up and held: structurally the same \textit{eidos} that perception first received, now held as the illuminated presence in which the thing itself can be seen and thought.

Because \textit{phantasia} is this illuminated holding-open of the \textit{eidos}, it grounds rather than merely accompanies the higher cognitive operations Heidegger names: thinking (``the matters are there in the making-present'' ), memory (recollection as a specific kind of making-present in which the ``having-once-experienced'' is also known as such ), and deliberation (``the thorough deliberating of a matter, when I do not have it perceptually present'' ). This is the deeper sense of Aristotle's claim that ``the soul never thinks without an image [\textit{phantasma}]'' (\textit{On the Soul} III.7, 431a16–17): not a contingent psychological condition on intellection, but the indication that \textit{phantasia} is the bodily ground of \textit{noein} itself---``if, however, even \gk{νοεῖν} {\textit{noein}} [the thorough deliberating of a matter, when I do not have it perceptually present] is something like a \gk{φαντασία} {\textit{phantasia}} or cannot be without \gk{φαντασία} [\textit{phantasia}], then thinking too could not be without standing in the context of the entire life of a human being'' (134).\footnote{I have placed the transliteration of the Greek words used by Heidegger within {} to avoid confusion, as this quotation already contained traditional editorial use of [brackets].} $A_2$ is precisely the moment at which the soul holds the \textit{phantasma} of the perceived particular under the description ``image-for-the-mind'': the moment, that is, at which every higher cognitive act of human being-there shows itself as bodily through \textit{phantasia}, since ``the entire being of human beings is determined so that it must be apprehended as the \textit{bodily being-in-the-world} of human beings'' (134; my emphasis).

Accordingly, in the \textit{phantasma} we see what Heidegger calls the ``making-present'' function of \textit{phantasia}, or what I will the \textit{taking-as-something}. Where bare \textit{aisthēsis} reveals proper sensibles only (white and spherical) the \textit{phantasma} is \textit{taken-as-baseball}.\footnote{Nussbaum's reading of \textit{phantasia} as the faculty that ``ties abstract thought to concrete perceptible objects or situations'' (Nussbaum 265) supplies the canonical scholarly anchor for the \textit{taking-as-something} reading the present section deploys. Hawhee extends the Nussbaumian \textit{seeing-as} into a rhetorical-vision frame: the \textit{phantasma} bears not just cognitive content but the affective-evaluative valence that makes the object \textit{matter as} what it is taken as (Hawhee, ``Looking Into Aristotle's Eyes: Toward a Theory of Rhetorical Vision,'' \textit{Advances in the History of Rhetoric} 14.2, 2011, esp.~pp.~147--155). The \textit{taking-as-baseball} case the section develops is, on these readings, not a cognitive add-on to a prior neutral perception but the structural mode in which the \textit{phantasma} discloses its content at all.}\footnote{NOTE TO SYSTEM: include at least a brief quotation from Hawhee. the current page range is too broad without actually mentioning something specific and implies that Hawhee is deliberately expanding Nussbaum's reading into her own concept. This must actually be the case, Hawhee must actually cite nussbaum and develop her thinking in order to make this claim.} To further clarify this function, we can turn to one of the most important Aristotelian observations this section deploys throughout: ``I want to drink says appetite; this is drink, says sense or imagination or thought: straightaway I drink'' (\textit{Movement of Animals} 7, 701a32-33). In this instance of the most basic appetitive action, Aristotle notes three different faculties which can impel a living creature to move: sense (\textit{aisthesis}), imagination (\textit{phantasia}), or thought (\textit{nous}).\footnote{Caston's reading of why Aristotle introduces \textit{phantasia} as a distinct faculty supplies the structural rationale for the present section's three-case treatment of \textit{MA} 7, 701a32--33: ``Phantasia thus stands midway between these different powers. It is more basic than the intellectual powers, extending throughout the animal kingdom, while also being capable of error. In this way, it is ideally suited to explain the errors that are manifest in the behavior of even nonrational animals'' (Caston 45). \textit{Phantasia}'s placement between aisthēsis and nous is what makes the three impetuses Aristotle catalogs at \textit{MA} 7 structurally homologous: each can substitute for the other as the cognitive route to the same action because each operates on (or as) the same \textit{phantasma}-borne aspectual presentation.} 

Let's begin with the first case. For example: You are thirsty after returning home from a run and walk into your kitchen. You see a cool glass of water on the counter. The very \textit{taking-as-something} (taking-as-drink) indicates a \textit{phantasma} is deployed: the \textit{eidos} ``drink'' is deployed \textit{with} the live visual data delivering the recognition as ``drink.'' In this moment, the proper sensibles are actualized: visual (color, transparency, cool-looking condensation on the glass). At the visual channel, as the sensible object is actually present, \textit{epithymia} is co-given: the clear, cool-look of condensation is intrinsically pleasant. However, \textit{resonant epithymia} also plays a role as a \textit{dunamis}, as a potentiality. The \textit{phantasma}-borne hedonic tonality of ``drink-as-realizable-good-to-be-pursued'' is applied to \textit{this} perceived particular; the drinking-engagement is \textit{anticipated}, but not yet \textit{realized}. You then grab the glass and drink straightaway. The \textit{eidos} ``drink'' is \textit{being realized} in the engagement and the proper-sensibles are actualized: taste (sweet, bland, mineral-like) and tactile (cool against the lip, wet and cold against the throat). In the moment of drinking, \textit{epithymia} is present and active where before they were merely anticipated. \textit{Resonant epithymia}, ``drink-as-realizable-good-to-be-pursued'' is actualized; the same \textit{phantasma}-borne hedonic tonality of ``drink-as-good'' is \textit{no longer anticipatory} but is being \textit{actualized} in the drinking. The \textit{resonant epithymia} has not disappeared or been replaced with \textit{epithymia}, it has simply shifted from \textit{dunamis} (potentiality) to \textit{energeia} (actuality). 

For the second case, imagination or \textit{phantasia} ``says'' this is drink. Let's say you are hiking in mid-summer, thirsty, and there is no water in sight. You have a canteen in your pack. You stop and an image arises: the canteen filled this morning, its weight, the cool water inside. This is the clearest case of \textit{phantasia} making-present something that is absent through the \textit{phantasma}. In this instance, we have the \textit{eidos} ``drink'' and the more specific \textit{phantasma} of the canteen-with-water present without any active sensory data of the canteen to engage with; no proper sensibles are actualized. There is no \textit{pathos-simpliciter} beyond the negative hedonic tonalities present which give rise to the drink-as-realizable-good: hot environment, dry throat, physical fatigue, etc. The \textit{phantasma}-borne hedonic tonality of drink-as-good-from-memory, drink-as-realizable-good, is held entirely \textit{without} current proper-sensibles of drink. This is \textit{resonant epithymia} in its purest mode: hedonic tonality riding on the \textit{phantasma} alone. The \textit{eidos} ``drink'' carries the goodness, and the \textit{phantasma} of canteen-with-water-in-pack specifies the realization-path; the \textit{resonant epithymia} motivates the action of reaching inside the pack for the canteen. The decisive contrast with the first case is that the \textit{phantasma}-affect (\textit{resonant epithymia}) is layered with drink-\textit{epithymia} (visual and/or tactile) because a drink-proper-sensible is firing. In the second case, the \textit{phantasma}-affect is \textit{not} layered with drink-\textit{epithymia}, because no drink-proper-sensible is firing. The proper-sensible that \textit{are} firing are those of the \textit{thirst-state} and being-in-that-specific-environment, not of drink---and their hedonic tonality (heat-pain, throat-distress) is exactly what makes the \textit{resonant epithymia} on the drink-\textit{phantasma} so motivating. 

The last case, thought or \textit{nous} says ``this is drink'' is purely inferential and categorical. Once again, you are hiking in unfamiliar terrain and are thirsty but have no water. However, it has rained very recently and you begin looking for areas that could still contain rainwater. You come over a rise a see clear liquid pooled in a basin-like depression in the rock. You have no specific memory-image of this exact liquid-in-basin. \textit{Nous} runs deliberative inference synthesizing various \textit{phantasmata}---recent rain, rain water, no contaminants, clean rock-basin---and concludes this-is-drink. In this case, the \textit{eidos} ``drink'' is held under universal premise (water-as-drink, drink-as-life-sustaining, drink-as-good); and, there is no specific \textit{phantasma} or direct sensory perception of any particular drink-source. \textit{Nous} operates on the \textit{phantasma} in its most generic, abstract mode. In the first half of this case, there is no proper sensibles or \textit{epithymia} for ``drink''; only environmental proper sensibles and their accompanying \textit{epithymia} and \textit{resonant epithymia}: drink-as-good-held-universally. The \textit{phantasma}-borne hedonic tonality is at its most abstract. Then a transition occurs when we see the puddle in the rock-basin: the abstract universal \textit{phantasma} is now applied to the perceived particular via the inference. Once the inference completes, the \textit{phantasma} is concretely integrated with the current visual data, just as in the first case of \textit{aisthesis}. We now have similar proper sensibles, \textit{pathos simipliciter} in the visual channel, and the abstract \textit{resonant epithymia} is now concretized: drink-as-good (held universally) becomes drink-as-realizable-good (applied to the particular). Once you kneel and take a drink, we have returned to the first case and \textit{epithymia} activates and the \textit{resonant epithymia} transitions from realizable to being-realized, from potential to actual. 

What we see through all three cases Aristotle identifies is that \textit{resonant epithymia} on the \textit{phantasma} ``drink'' is present everywhere the \textit{eidos} is held: presently deployed (when seeing the drink-as-realizable-good), retained from memory (memory of the canteen in our pack, memory of recent rain), prospectively applied (search for places which could hold rainwater), and even atemporally (\textit{nous} holds the \textit{phantasma} abstractly in drink-as-good-universally). 

The three cases just considered isolated the structural fact that the \textit{phantasma}'s \textit{taking-as-something} can be activated by any of the three discriminative faculties Aristotle identifies---by \textit{aisthēsis} when the object is present, by \textit{phantasia} when it is absent (whether recalled from memory or anticipated), by \textit{nous} when it is held universally---and that each such activation, when conjoined with appetite, can issue ``straightaway'' in action. But that is only one of the trajectories the \textit{phantasma} admits. Rather than being discharged into immediate action, the \textit{phantasma} can equally be \textit{held} by the soul and engaged with cognitively; and the modes of \textit{phantasma}-deployment that the second and third cases brought to light---retrospective, prospective, and universal-abstract---reappear here as the temporal orientations of three distinct intentional directions the soul can take toward the \textit{phantasma} it holds. Three orientational modes must accordingly be distinguished: \textit{intellection} (\gk{νοῦς}), \textit{memory} (\gk{μνήμη}), and the \textit{discursive} mode that Aristotle designates \textit{deliberative phantasia} (\gk{φαντασία βουλευτική})\footnote{Aristotle introduces \textit{phantasia bouleutik\=e} at \textit{De Anima} III.11, 434a5--10 and develops its distinguishing capacity at 434a16--21: the deliberative grade ``exists in those animals which can calculate; for whether one will do this or that is already a work of calculation; and there must be a single measure, for one pursues what is greater. Hence one has the ability to make one \textit{phantasma} out of many \textit{phantasmata}'' (434a7--10; \gk{ἓν ἐκ πλειόνων φαντασμάτων ποιεῖν}). The combination of multiple \textit{phantasmata} into a single object of cognitive address is the defining capacity that distinguishes rational engagement with the image from the non-rational animal's engagement with a single present image.} and which divides, as will be shown, into practical and speculative (theoretical) sub-domains. The three share a common structural character: each is a \textit{non-committal} cognitive engagement with the \textit{phantasma}, in the precise sense that none of them, taken on its own, yet makes a truth-claim about what the \textit{phantasma} presents. They differ from one another in the intentional direction each brings to the engagement and in the temporal orientation each presupposes: intellection holds the \textit{eidos} atemporally, as universal; memory deploys the \textit{phantasma} retrospectively, as past; the discursive mode projects it prospectively, whether toward future action in the practical case or toward inferred truth in the speculative. Alongside these three orientational modes, a structurally distinct dimension enters the engagement: \textit{doxa}, the propositional \textit{taking-as-something-true} through which the rational animal commits the soul to a truth-claim about what the \textit{phantasma} has been \textit{taken-as}. The present section treats the three orientational modes in turn; the \textit{doxastic} dimension will be treated last, where it will be shown to operate not as a fourth parallel mode but as an \textit{orthogonal committal layer} that can engage with any of the three orientations.

The \textit{phantasma}, then, can be taken up under several distinct intentional directions, each issuing in a different sort of completed cognition. Three orientational modes must be distinguished: intellection (\textit{nous}), memory, and that mode which Aristotle designates deliberative \textit{phantasia} and which divides, as will be shown, into practical and speculative or theoretical sub-domains. These three modes share a common structural character as non-committal cognitive engagements with the \textit{phantasma}, and they differ from one another in the intentional direction each brings to the engagement and in the temporal orientation each presupposes. Alongside these three orientational modes, a structurally distinct committal dimension enters the engagement: \textit{doxa}, the propositional \textit{taking-as-true-or-false} through which the rational animal commits the soul to a truth-claim about what the \textit{phantasma} is \textit{taken-as}. The present section treats the three orientational modes in turn; the \textit{doxastic} dimension will be treated last in this section, where it will be shown to operate not as a fourth parallel mode but as an orthogonal committal layer that can engage with any of the three orientations (\textit{De Anima} III.3, 428a19--24 on \textit{doxa} as entailing \textit{pistis} where \textit{phantasia} does not).

\inlinenote{One note: the noein/phantasia citation now carries two kinds of brackets — your editorial transliterations [\textit{noein}], [\textit{phantasia}] and the translator's own bracketed gloss [the thorough deliberating of a matter, when I do not have it perceptually present]. Both conform to standard practice (square brackets = editorial insertion), but a reader cannot tell at a glance which inserter is which. If that matters for your audience, the cleanest disambiguation is a single footnote at first occurrence: "Greek terms in square brackets are my transliterations; longer bracketed material is from the Rojcewicz/Schuwer translation."}

\subsection{Intellection: The \textit{Phantasma} as Bearer of the Universal}

The first orientational mode is intellection, in which the \textit{phantasma} is addressed as bearing a universal. Aristotle insists that ``the soul never thinks without an image'' (\textit{On the Soul} III.7, 431b2) and that ``the thinking faculty thinks the forms in the images'' (431b12-13). Intellection (\textit{nous}), then, operates upon the \textit{phantasma} by abstracting from its particular determinations the universal form that it instantiates; the \textit{phantasma}, in other words, serves as the substrate from which intellect extracts intelligible content. Accordingly, while \textit{phantasia} ``is different from either perceiving or discursive thinking,'' Aristotle clarifies, ``it is not found without sensation, or judgement without it'' (\textit{On the Soul} III.3, 14-15). However, what \textit{nous} grasps in the \textit{phantasma} is not the particular image as particular but what is universal within it. The intelligible form is abstracted from the sensory image and apprehended as holding across all the instances of which this particular \textit{phantasma} is one. Thus, the rational animal, having intellected the \textit{phantasma} of this man as bearing the universal \textit{human}, thereby grasps something that applies not only to the man presently imaged but to every man that has ever or will ever exist.\footnote{White's reading of \textit{De Anima} III.7--8 further articulates how the \textit{phantasma} can come to bear a universal once abstracted through intellection: \textit{phantasia} is the storehouse of forms that ``serves as a step in man's progression beyond time altogether, enabling him at certain points and for certain periods to pass beyond the confinements of sense to a consideration of the timeless nature of things'' (White 503). Intellection abstracts the universal that \textit{phantasia}'s preservation has already held in availability.} Thinking admits both falsehood and truth, and it exists only where there is rational discourse; thus the \textit{phantasma}, insofar as it serves as the bearer of the universal for \textit{nous}, must be a content that is articulable in \textit{logos}.\footnote{Heidegger develops the \textit{logos}/\textit{path\=e} relation at \textit{Basic Concepts of Aristotelian Philosophy} along three load-bearing axes: rhetoric is ``the interpretation of being-there with regard to the basic possibility of speaking-with-one-another'' (\textit{BCAP}, p.~95); deliberative oratory turns on three factors --- ``(1) \textit{ethos}, the `comportment' of the speaker; (2) \textit{pathos}, the `disposition' of the hearers; (3) the matter itself'' (\textit{BCAP}, p.~110); and the master claim follows at p.~176: ``the \gk{πάθη} are not merely an annex of psychical processes, but are rather \textit{the ground out of which speaking arises, and which what is expressed grows back into}.'' \textit{Phantasia} is what makes the universal sayable; \textit{logos} is what makes the said universal articulately bear truth or falsehood. The \textit{phantasma}-bearing-universal is the condition for both.} The soul does not merely gaze upon the image but speaks through it, formulating discourse that articulates the universal in conceptual form that the \textit{phantasma} bears.\footnote{The ``universal holds across instances'' claim is grounded at \textit{Metaphysics} Z.10 (1035b27--1036a25), where Aristotle distinguishes the universal-as-form (\textit{logos}) from the particular composite. Caston gives the corresponding cognitive-architecture reading: ``by making \textit{phantasia} (rather than belief) the common currency among mental states, Aristotle emphasizes a form of intentionality more basic than the conceptual'' (Caston 1996, \textit{Why Aristotle Needs Imagination}, p.~47). The \textit{phantasma} thus serves as the pre-conceptual vehicle of intentional content; \textit{nous}'s abstractive operation extracts the universal that the \textit{phantasma} bears without abstracting from \textit{phantasia} itself.} 

The temporal orientation of intellection is, accordingly, \textit{atemporal}. The universal, once grasped, is not indexed to past, present, or future; it holds across all the times at which its instances might occur and is not itself situated within the temporal continuum that conditions the other two modes. This atemporality distinguishes intellection structurally from memory, whose object is constitutively past, and from the deliberative mode, whose objects are either projected future possibilities or contemplated unchangeable truths. Intellection alone operates beyond the temporal horizon, for the universal that \textit{nous} grasps has no date. However, as Heidegger is right to note, even the \textit{nous} ``of human beings is not \textit{pure}'' (136; my emphasis) but is grounded, bodily, in \textit{phantasia}, is determined as bodily being-in-the-world. Heidegger's gloss adds an implicit temporal dimension: the atemporality of intellection is the atemporality of \textit{eidos}-grasping. The universal has no date, but its grasp is conditioned on the bodily, temporal availability of the \textit{phantasma}through which it shows. This atemporality also marks the structural boundary between intellection and \textit{doxa}. To intellect the universal \textit{human} is to grasp it as a conceptual content; to hold a \textit{doxa} that ``this man is human'' is to propositionally ratify a truth-claim in which the universal figures. The grasping and the ratifying are distinct operations. \textit{Nous} abstracts; \textit{doxa} asserts. \textit{Nous} can operate without \textit{doxa}---one can grasp what \textit{human} is without having yet formed a belief about any particular human---and \textit{doxa} can operate on the universals that \textit{nous} supplies. The $A_3$ actuality produced by intellection is thus the universal grasped; whether a \textit{doxa} about that universal also arises is a further question to be taken up by the committal dimension in the section that follows.

\subsection{Memory: The \textit{Phantasma} as Likeness of the Past}

The second orientational mode is memory, in which the \textit{phantasma} is addressed as a likeness of something past. Memory, as developed in \textit{On Memory}, presents a distinct cognitive orientation upon the \textit{phantasma}, one that treats the image not as the bearer of a universal but as a likeness of a past perceptual encounter: ``There is no such thing as memory of the present while present,'' Aristotle writes, ``for the present is object only of perception, and the future, of expectation, but the object of memory is past'' (449b25-27). The object of memory is definitionally, as Frede observes, ``a past experience \textit{qua} past'' (285). The ``\textit{qua} past'' is essential, and it marks the distinction between memory and \textit{phantasia} more generally. Only beings that can perceive time can remember, ``and the organ whereby they perceive time,'' as Aristotle states, ``is also that whereby they remember'' (\textit{On Memory} 1, 449b28-30). Thus, the \textit{phantasma} is functioning not merely as a spatial representation of an absent object but a temporally indexed image that carries within itself the mark of pastness. As such, not every \textit{phantasma}is a memory. An image of an apple may arise in the mind as a free-floating representation: a \textit{phantasma} of apples in general, unassociated with any particular past encounter. In addition to the \textit{phantasma}, memory requires the recognition that the \textit{phantasma} is the likeness of something one has ``formerly heard or perceived or thought'' and is, Aristotle concludes, ``neither perception nor conception, but a state or affection of one of these, conditioned by lapse of time'' (\textit{On Memory} 1, 449b22-25). The lapse of time is not external to memory as its condition but internal to what memory discloses: the past is marked \textit{as past} and felt \textit{as past} within the memorial act itself. Heidegger's existential analytic casts additional light upon this structure. The ecstatic temporality that he describes as \textit{Zeitlichkeit}---unfolding in the three ecstases of future, having-been, and present---provides a framework within which the temporal indexing of the memory-\textit{phantasma} can be understood not as a mechanical timestamp but as a mode of Dasein's self-relation. The having-been (\textit{Gewesenheit}) is not a dead past but a living dimension of Dasein's being that continues to determine its present understanding and future projection (\textit{BT} 449-451). Aristotle's memory-\textit{phantasma}, similarly, is not a dead copy of a vanished experience but a living \textit{resonant kinēsis} that continues to orient the soul's cognitive engagement toward the past from which it derives.\footnote{Heidegger develops ecstatic temporality (\textit{ekstatische Zeitlichkeit}) at \textit{Being and Time} §65 (H.323--331) as the unitary structure of \textit{Zeitlichkeit}: ``Temporality temporalizes itself \ldots\ in the unity of ecstatic future, having-been, and Present'' (H.328--329). Each ecstasis is a ``rapture toward'' (\textit{Entr\"uckung zu}) a horizonal schema; the three together constitute the existential meaning of care (\textit{Sorge}). The having-been ecstasis (\textit{Gewesenheit}) is the existential meaning of disposedness (\textit{Befindlichkeit}); it is distinguished from past-as-no-longer (\textit{Vergangenheit}) as ``a `positive' ecstasis'' rather than mere absence: ``Dasein, as long as it is, is constantly its own `not-yet,''' and the having-been remains operative as the existential ground of present comportment (SZ §65, H.339--340). The full historicality unfolds at §§72--77 (H.372--404), where the unity is reconstructed as Dasein's stretching-along (\textit{Erstreckung}) between thrown beginning and projected end. The memory-\textit{phantasma}, on the present analysis, is the Aristotelian correlate of this \textit{Gewesenheit} structure: not a copy of a vanished episode but a living orientation that conditions the soul's current cognitive engagement.}

Heideger's analysis of recollection in \textit{Being and Time} (section 72-77) (anticipated in his discourse on Aristotle's \textit{hexis}\footnote{``The frequently is, precisely, that which characterizes the temporality of being-there. Aristotle cannot say \gk{ἀεί} (\textit{aei}) insofar as human being-there does not so comport itself constantly and always. It can constantly be otherwise. The \textit{always} of a being like being-there is the \textit{frequently of repetition}. It is the being-there of human beings, as determined by \textit{historicality}, to see entirely different time-contexts in relation to which the remaining time-determinations break down'' (129).} \hl{(Repeat quote in the final settled disposition section.)} provides a temporal framework: memory is a mode of \textit{Gewesenheit}, a living dimension of \textit{Dasein}'s being that continues to determine present understanding and future projection. The having-been is not a dead past but a living orientation that conditions the soul's current cognitive engagement. 

The temporal orientation of memory is, therefore, \textit{retrospective}.\footnote{Struever's reading of \textit{Alltäglichkeit} as ``timefulness'' --- the radical specificity of time, the care for tense that defines practical life --- supplies the matched political-philosophical reading of memory's retrospective character: ``timefulness colonizes \textit{Being and Time}'s temporal program from within'' precisely because each \textit{pathos}-actualization is shaped by the having-been of prior actualizations, which the memory-\textit{phantasma} carries forward as the dispositional ground of present engagement (Struever 117). The memory-\textit{phantasma}'s retrospective character is not merely a cognitive timestamp but the operative form of the \textit{Gewesenheit}-ecstasis at the level of the rational soul's discriminative engagement with the \textit{phantasma}.} The \textit{phantasma} is addressed not as bearing a universal and not as composable into a future possibility, but as point backward toward the past encounter of which it is a likeness. The $A_3$ actuality produced by memory is the past recalled and recognized as past---a content indexed to the temporal horizon of its original occurrence and apprehended from the position of the present as belonging to that prior horizon. Accordingly, memory furnishes the materials upon which the other modes may subsequently operate. A memory may provide the particular instance from which intellection abstracts a universal; a repertoire of memories may provide the \textit{phantasmatic} materials that the deliberative mode combines into future projections; a settled memory memory may function, like a settled \textit{doxa}, as an additional unmoved originator that shapes how a current \textit{phantasma} is taken up in subsequent engagements. Memory is, in this sense, the temporal deposit upon which the rational animal's cognitive life draws. 

\subsection{The Discursive Mode: Deliberative and Speculative Thinking}

The third and final orientational mode differs from the other two in a structural respect: it operates not upon a single \textit{phantasma} but upon several, synthesizing them into a unified object of thought that none of the constituent \textit{phantasma} alone could furnish. Aristotle specifies the synthesizing character in the passage that grounds his analysis of practical deliberation: ``whether this or that shall be enacted is already a task requiring calculation; and there must be a single standard to measure by, for that is pursued which is greater. It follows that what acts in this way must be able to make a unity out of several images'' (\textit{DA} III.11, 434a5--10). The defining capacity of this mode is the combining of multiple \textit{phantasmata} into a single object of cognitive address, and it is this synthetic capacity that distinguishes rational engagement with the \textit{phantasma} from the non-rational animal's engagement with a single present image. Within this discursive, synthetic mode two domains must be distinguished: a practical domain (deliberative or anticipatory \textit{phantasia}) and a theoretical domain (speculative or theoretical thinking). Both domains share the underlying synthetic operation upon multiple \textit{phantasmata}; they differ in the nature of the object they synthesize toward and in whether the synthesis issues in action or terminates in contemplation. 

Where sensory \textit{phantasia} generates \textit{phantasmata} that are direct residues of past perceptions, deliberative \textit{phantasia} synthesizes multiple \textit{phantasmata} into composite images representing states of affairs that have not yet happened and may never happen. It operates to project unrealized possibilities and to measure them against one another for the sake of action. Two interlocking sub-functions constitute this practical operation. The \textit{projective-synthetic} sub-function recombines \textit{phantasmata} drawn from the available repertoire into images of states not yet actual: possible futures, anticipated outcomes, unrealized configurations. This is the productive capacity that goes beyond mere retention; the phantastic faculty here constructs rather than merely preserves, assembling from the materials of past experience images of what might yet come to pass. The \textit{calculative-practical} sub-function then measures the projected possibilities against one another by a single standard to determine which course of action shall be enacted. It is here that deliberative \textit{phantasia} converges with practical thought: the projected image is evaluated not merely as possible but as to-be-pursued or to-be-avoided, and the deliberative process terminates in a determination of the greater good.\footnote{Gonzalez's reading of \textit{Rhetoric} III.1.1404a11 supplies the rhetorical articulation of this synthetic-evaluative operation: \textit{phantasia}'s combinatorial work at the rhetorical level operates through \textit{lexis} (style) as the medium in which the audience's \textit{phantasmata} are reshaped and given the \textit{pistis} (conviction-bearing weight) that \textit{doxa} requires (Gonzalez 99-131; esp. 113-119 on  \textit{phantasia} \textit{pistis}). The structural similarity is clear: at the level of individual cognition, deliberative \textit{phantasia} synthesizes multiple \textit{phantasmata} into a single calculable course of action; at the rhetorical level, the orator's \textit{lexis} synthesizes the audience's available \textit{phantasmata} into a single \textit{doxa}-bearing presentation.} (NOTE TO SYSTEM: I need at least a brief quote from Gonalez. There might also be an opportunity to include O'Gorman as he talks about this rhetorical element specifically as well.)

Therefore, deliberative \textit{phantasia} is essentially prospective or future-oriented; it generates the image of a prospective good or evil that serves as the object of desire and the terminus of practical reasoning. The distinction between sensory and deliberative \textit{phantasia} is not merely a difference of degree but a difference of cognitive orientation: sensory \textit{phantasia} preserves, whereas deliberative \textit{phantasia} projects and constructs. This dual structure---synthetic projection plus calculative evaluation---is what distinguishes rational action from appetitve reaction. The voluntariness of deliberation \textit{phantasia} is a crucial element differentiating it from its sensory counterpart: whereas sensory \textit{phantasmata} arise unbidden as residues of perception, deliberative \textit{phantasmata} are to some degree under the control of the rational agent, who can summon, combine, and arrange images in the service of practical reasoning. As such, deliberative \textit{phantasia} allows rational souls to hold back from what is immediately present by projecting what is future, and this I holding-back requires---indeed presupposes---a sense of time: ``while thought bids us hold back because of what is future, desire is influenced by what is just at hand: a pleasant object which is just at hand presents itself as both pleasant and good, without condition in either case, because of want of foresight into what is farther away in time'' (\textit{On the Soul}, III.10, 433b5-10). The sense of time is thus co-implicated with deliberative \textit{phantasia}: the capacity to calculate and the capacity to perceive time cannot, on Aristotle's account, be separated without destroying both. Additionally, deliberation, as Heidegger parses it, is a ``seeking,'' ``a being-after a definite view that I want to achieve. Through deliberating, I want primarily to come to the \gk{τέλος} (\textit{telos}) of a \gk{δόξα} (\textit{doxa}); it is not a yes. I only want to appropriate a definite view regarding a set of facts'' (96).\footnote{Kisiel's reading of Heidegger's Aristotelian lectures specifies the political-philosophical stakes of this seeking: deliberation is grounded in the cultivated attunement (\textit{Stimmung}) of the auditor by the skilled speaker, ``so that the attentive listener will then speak with the speaker by speaking after him or her (\textit{Nachreden}) in contagious attunement, suggests that speech finds its deepest roots in mood'' (Kisiel 145, citing BCAP, p.~177). The deliberating soul's seeking-after-a-definite-view is therefore not isolatable from the social-rhetorical context in which the deliberation unfolds: rational deliberation is constitutively informed by the dispositional ground (cultivated \textit{hexeis} and active \textit{pathē}) within which alone the relevant \textit{doxai} can come into view.}

Speculative thinking constitutes the second dimension of the discursive mode, operating upon the \textit{phantasma} not in the service of practical deliberation but for the sake of contemplation itself. Where deliberative \textit{phantasia} combines \textit{phantasmata} to construct images of prospective goods or evils, speculative thinking combines and recombines \textit{phantasmata} to explore intelligible structures like motion and time and the relation between them. What distinguishes speculative thinking from its deliberative counterpart is, then, the domain of its object. Deliberative \textit{phantasia} synthesizes toward the realizable---toward what might yet come to pass and is therefore measurable by the standard of the greater good. Speculative thinking synthesizes toward the unchangeable---toward what is and cannot be otherwise and is accordingly not subject to practical measurement at all. Aristotle marks this distinction quite clearly: ``thought which is concerned with action'' differs from ``thought which is speculative and not concerned with action''; the former ``calculates means to an end,'' while the latter ``moves nothing'' (\textit{On the Soul} III.10, 433a13-15).  Practical though moves because its object is a realizable good that engages \textit{orexis}; speculative thought does not move because its object is not a realizable good but a truth that is what it is independently of anything the soul might do. However, here I must clarify that it does not move the animal in the sense of final action as there is no object of pursuit or avoidance; rather, speculative \textit{phantasia} is a movement insofar as it is a thinking-through-\textit{phantasmata} that discovers intelligible relations in the very process of its movement. The temporal orientation of speculative \textit{phantasia} is accordingly \textit{atemporal-contemplative}: its object neither lies in the past (as with memory) nor in the projected future (as with deliberative \textit{phantasia}) but stands outside the practical temporal horizon altogether, as what holds and cannot be otherwise. 

The $A_3$ actualities produced by the synthetic mode are, then, twofold. In the practical case, $A_3$ is the anticipated possibility that has been constructed from a synthesis of \textit{phantasmata}, evaluated against a single standard of measure, and either adopted or rejected as a course of action to be pursued. In the speculative case. $A_3$ is the truth contemplated---the universal, essence, or necessary connection grasped through the composed \textit{phantasmatic} materials---without issuing in any realizable good. The chain, then, terminates at $A_3$ for the speculative case: the contemplated truth moves nothing, and no \textit{orektikon} motion is initiated at $M_3 \to A_4$ on its basis. Practical deliberation, by contrast, does initiate such motion, for the deliberated possibility discloses, at $A_3$, a realizable good toward which \textit{orexis} can be moved. 

\footnote{Papachristou (2013) argues that Aristotle in \textit{De Anima} III distinguishes \textit{three}, not two, kinds of \textit{phantasia}: indeterminate (\gk{ἀόριστος φαντασία}) in imperfect animals (zoophytes, molluscs) that possess only the contact sense (\textit{De Anima} III.11, 433b31--434a5); sensitive (\gk{αἰσθητική φαντασία}) in non-rational animals with multiple senses; and calculative or deliberative (\gk{λογιστικὴ φαντασία} or \gk{φαντασία βουλευτική}) in rational beings (\textit{De Anima} III.10, 433b29; III.11, 434a5--10). The textus receptus is Thomas Aquinas, \textit{Sentencia libri De anima}, Liber II: ``\textit{Dicendum est igitur, quod animalia imperfecta, ut in tertio dicetur, habent quidem phantasiam, sed indeterminatam, quia scilicet motus phantasiae non remanet in eis post apprehensionem sensus: in animalibus autem perfectis remanet motus phantasiae, etiam abeuntibus sensibilibus}'' (Aquinas, \textit{Sent.\ De an.} II, in Papachristou 2013, p.~16, fn.~57; trans.\ Kocourek: ``It must be said, therefore, that imperfect animals, as is said in the third book, do really have phantasy, but it is one which is indeterminate because the motion of phantasy does not remain in them after the apprehension of the sense; however in perfect animals the motion of phantasy remains even after the sensible thing is gone''). Papachristou's tripartite scheme is the textual ground for the present chapter's distinction of orientational modes operative within the rational-animal case; her three-grade scheme operates at a different axis (animal-class-stratification) from the three-orientational-modes scheme deployed here (within-rational-animal modes), but both depend on the same Aquinas/Papachristou anchoring of \textit{phantasia bouleutik\=e} as the discriminating capacity of rational beings.}

\subsection{The Unified Cognitive Apparatus}

These three orientational modes---intellection, memory, and the discursive mode in its practical and speculative domains---should not be mistaken for three separate faculties operating independently within the rational soul. They are, rather, analytical descriptions of a single rational cognitive apparatus operating upon the \textit{phantasma} under distinct intentional directions. They share a common structural role within $M_2 \rightarrow A_3$: in each mode, the unmoved originator is $A_2$ (the \textit{phantasma}, or in the synthetic case, the repertoire of \textit{phantasmata} available to the soul); the moved mover is the rational cognitive apparatus in its respective orientation; that which is moved is the rational animal, cognitively transformed by the engagement. Each orientation takes up the same formally determined image and engages it under a different aspect: intellection abstracts the universal; memory indexes the image to the past; deliberative \textit{phantasia} projects it into an imagined future; speculative thinking traces its intelligible relations for the sake of contemplation. They also differ in the temporal orientation that each direction presupposes: atemporal for intellection, retrospective for memory, prospective for deliberative \textit{phantasia}, atemporal-contemplative for speculative \textit{phantasia}. Thus, the \textit{phantasma} functions as the hinge upon which the entire cognitive life of the rational soul turns, and the passage from \textit{resonant kinēsis} to determinate image---is the enabling condition of all higher cognition. What I have been calling ``orientation modes'' we could also call, in a Heideggerian register, different modes of disclosure; what I speak of ``temporal orientations'' Heidegger speaks of ``ecstases of \textit{Zeitlichkeit}.''\footnote{The three ecstases (\textit{Ekstasen}) of temporality --- futural (\textit{Zukunft}), having-been (\textit{Gewesenheit}), Present (\textit{Gegenwart}) --- are introduced at \textit{Being and Time} §65 (H.323--331) as the unitary structure of \textit{Zeitlichkeit}; each ecstasis is a ``rapture toward'' (\textit{Entr\"uckung zu}) a horizonal schema, and the three together constitute the existential meaning of care (\textit{Sorge}). The futural ecstasis's horizonal schema is the ``for-the-sake-of-itself,'' having-been's is ``that to which thrown,'' the Present's is the ``in-order-to'' (SZ §69c, H.365). The three orientational modes of the present chapter --- intellection (atemporal universal), memory (retrospective), discursive (prospective practical / atemporal-contemplative speculative) --- map onto, but are not identical with, Heidegger's three ecstases: the chapter's modes are the Aristotelian cognitive operations through which the soul takes up the \textit{phantasma}; Heidegger's ecstases are the structural-temporal raptures within which any such taking-up is possible at all.}\footnote{The three-factor kinetic schema is given at \textit{De Anima} III.10, 433b13--18: ``these three are required: that which moves, the means by which it moves, and that which is moved'' (433b13--16). At the perceptual stage (§1.2) the unmoved originator was the sensible object exercising its proper quality; at the \textit{phantastic} stage ($M_1 \rightarrow A_2$) the unmoved originator is the resonant trace; at the cognitive stage ($M_2 \rightarrow A_3$) developed here, the unmoved originator is $A_2$ (the \textit{phantasma} stabilized at the central organ); the moved mover is the rational cognitive apparatus in its respective orientation; the moved is the rational animal cognitively transformed.} The atemporal universal, the past as past, the future possibility, and the unchangeable truth are four ways in which something can be made present without being sensorily present. 

These modes may operate simultaneously or in succession; the same \textit{phantasma} can simultaneously bear a universal for intellection, a temporal index for memory, a deliberative projection for practical reasoning, and a formal structure for speculative contemplation, and these multiple engagements interpenetrate in the concrete cognitive life of the rational soul. These modes are not mutually exclusive orientations but distinguishable axes of engagement that the rational soul may occupy in any combination. Each mode terminates in a distinct sort of $A_3$ content: the universal grasped (intellection), the past recognized as past (memory), the possibility constructed and evaluated (deliberative), the truth contemplated (speculative). These are the completed cognitive actualities that $M_2 \to A_3$ produces when the rational cognitive apparatus engages the \textit{phantasma} in its orientational capacity. The \textit{phantasma}, therefore, is the bodily-temporal medium in which the soul's disclosedness is actualized. In each case, the content of $A_3$ can, for the rational animal, be taken up by a further operation: the orthogonal committal dimension of \textit{doxa}, which ratifies the content as true and binds the soul to the ratification involuntarily. It is to \textit{doxa}, and to its structural character as an orthogonal committal dimension rather than a fourth parallel mode, that the next section turns. 

\subsection{\textit{Doxa}: The Soul's Involuntary \textit{Taking-as-Something-True-or-False}}

\textit{Doxa} occupies a distinctive position in Aristotle's cognitive architecture because it is not a fourth orientation mode alongside intellection, memory, and discursive thought; bur rather an orthogonal dimension that attaches to any content produced by those orientations. Aristotle opens his contrast between \textit{phantasia} and \textit{doxa} with a claim about the will. ``Imagining,'' he observes, ``lies within our own power whenever we wish (e.g.\ we can call up a picture, as in the practice of mnemonics by the use of mental images), but in forming opinions we are not free: we cannot escape the alternative of falsehood or truth'' (\textit{DA} III.3, 427b17--21). The asymmetry is total. \textit{Phantasia} is \gk{ἐφ' ἡμῖν}---within our power, available to summoning, at the disposal of any soul that cares to practice the mnemonic art. \textit{Doxa} is \gk{οὐκ ἐφ' ἡμῖν}---not within our power, not available to decision, withheld from the will precisely because it is answerable to something beyond the will. 

Heidegger notes the difference between \textit{phantasia} and \textit{doxa} with his usual precision: 
\begin{adjustwidth}{0.5in}{0in}
    \gk{Φαντασία} [\textit{phantasia}]: the ``having-present'' of something without perceiving it directly, the mere ``presenting-itself.'' It can be true or false like \gk{δόξα} [\textit{doxa}]. It has both possibilities, but it has them, in a certain sense, only \textit{from without}, while \gk{δόξα} [\textit{doxa}] has the possibility \textit{in itself}. Within the sense of opining itself lies the ``can''---true or false. \gk{Δυνατόν---ἀδύνατον} [(]\textit{Dynatón---adýnaton/``possible''---``impossible''}]. (93-94)
\end{adjustwidth}
Furthermore, where the truthood or falsity of \textit{phantasia} ``is added to it,''\footnote{Heidegger develops the ``truth-from-without'' formulation at \textit{BCAP} pp.~93--94, in his analysis of \textit{doxa} (§15a): ``\textit{phantasia} is simple `having-present' without perceiving directly; \textit{doxa} has the possibility of truth or falsity \textit{in itself} (\textit{De Anima} III.3), while \textit{phantasia} is only [truth or falsity] from without; animals possess \textit{phantasia} too, but \textit{doxa} only belongs to beings with \textit{logos}'' (\textit{BCAP}, pp.~93--94). The structural distinction is that \textit{phantasia}'s truth-or-falsity is a property of the relation between the image and what it images, not a property of the image's own constitution; \textit{doxa}'s truth-or-falsity is constitutive of \textit{doxa} as such, because \textit{doxa} just is the assertoric uptake of the \textit{phantasma} under \textit{pistis} (\textit{De Anima} III.3, 428a19--24) --- the truth-claim is internal to \textit{doxa}'s own mode of being. Aristotle's source-claim is at \textit{De Anima} III.3, 428a3--4 (the \textit{phantasia}/\textit{doxa} contrast) + 428a18--21 (the \textit{pistis}-condition on \textit{doxa}).} sfor \textit{doxa} ``the being-able-to-be-true-or-false is already contained in the formation off the view'' (96). Where \textit{phantasia} is a simple making-present, taking-as-something, \textit{doxa} is the ``yes-saying'' that adds \textit{logos}-dependent ratification. Rationality is precisely what makes the \textit{taking-as-something-true-or-false} possible and why \textit{doxa} is only present in rational beings. One can imagine the coiled object in the shadowed corner of a shed as a snake; one cannot, having definitively seen that it is a rope, then decide to believe it is a snake. The capacity and the inability arise from a single source: \textit{doxa} is constitutively truth-or-falsehood-bearing, and it is answerable, accordingly, to how things \textit{are} rather than to how we might wish them to appear. \textit{Doxa} is the soul's involuntary taking-as-something-true or taking-as-something-false: when the \textit{phantasma} is engaged by intellection, memory, deliberation or speculation, \textit{doxa} pronounces upon the resulting content, committing the soul to the truth or falsity of what is presented. This committal is involuntary in the sense that it is not subject to the agent's deliberate control in the way that the summoning of deliberative \textit{phantasmata} is; rather, \textit{doxa} engages cognitive orientations. Thinking issues in propositions---articulations in \textit{logos}---and \textit{doxa} is the soul's stance toward those propositions as true or false; it is not a pre-propositional attitude but a response to propositionally structured content.

Where all higher cognition is grounded in the image-based, pre-propositional mode of being-in-the-world that \textit{phantasia} discloses, \textit{doxa}, as Heidegger illuminates, names the average everyday discoveredness of ``\textit{being-oriented in the world}, how human being-there \textit{initially has its world there in an average way}'' (\textit{BCAP} 101). \textit{Doxa}, as he reads Aristotle, is to ``have ``a view already'','' ``to be situated thus and so toward the matter'', ``a certain \textit{yes-saying} to that of which I have a view'' (93). Thus, Heidegger's description of \textit{doxa} reveals a positive ontological structure prior to and grounding both rhetorical persuasion and theoretical science the decisive feature of which being that \textit{doxa} is a ``\gk{φάσις}'' (\textit{phasis}), \textit{doxa asserts}: it \textit{holds} a view---as already in place---not as a result of inquiry but as the antecedent ground from which inquiry can even begin. 

Aristotle specifies the ontological distinctness of \textit{doxa} with three conditions that \textit{phantasia} characteristically lacks. ``Opinion,'' he writes, ``involves belief (for without belief in what we opine we cannot have an opinion), and in the brutes though we often find imagination we never find belief. Further, every opinion is accompanied by belief, belief by conviction, and conviction by discourse of reason'' (\textit{DA} III.3, 428a19--24). The three conditions are not redundant glosses upon a single requirement but a layered specification of what \textit{doxa} must possess in order to be what it is. Belief names the assent that distinguishes \textit{doxa} from the non-committal entertainment of an image: to opine is to take something as the case, not merely to present it to the mind. Conviction names the binding force of that assent: \textit{doxa} does not hold its content tentatively but is held by it, and it is this being-held that the soul cannot withdraw at will. Discourse of reason names the medium through which the assent and its binding force operate: \textit{doxa} is not the bare registration of an appearance but an assertoric articulation of that appearance under the categories that reason makes available.\footnote{Caston's reading of the \textit{phantasia}/\textit{doxa} contrast clarifies the structural asymmetry at issue: ``Phantasia has a wider extension than belief. Every animal that has belief also has phantasia, since belief is a kind of accepting or taking something to be the case ... and there is no acceptance without phantasia'' (Caston 45, citing DA 427b24--26). But the converse fails: ``It is not possible to believe just anything, but only what we judge to be right ... Phantasia, in contrast, allows us to conjure up images freely, independent of our views about what is true'' (Caston 45, citing DA 428a23--24, 427b17--20). The three-condition layered specification Aristotle gives at DA 428a19--24 (belief / conviction / discourse-of-reason) is therefore not a redundant gloss but the precise architectural specification of what \textit{doxa} must add to \textit{phantasia}'s prior \textit{taking-as} to qualify as truth-committal cognition.}

\textit{Doxa} and \textit{phantasia} can deliver contradictory content simultaneously, and this phenomenon demonstrates their ontological distinctness with clarity. Aristotle provides the canonical illustration: ``we imagine the sun to be a foot in diameter though we are convinced that it is larger than the inhabited part of the earth'' (\textit{DA} III.3, 428b2--4). The \textit{phantasma} is not corrected by the \textit{doxa}---it persists unchanged, presenting its testimony regardless of what reason concludes---but neither is \textit{doxa} compelled by the \textit{phantasma}'s apparent testimony. The soul holds two contents at once: the image of the one-foot sun, the belief in the supra-terrestrial sun, each fully present, neither yielding to the other. If \textit{doxa} were identical to \textit{phantasia}, or if it were merely a modification of the \textit{phantasma}'s content, the simultaneous holding would be impossible: the correction of the image would be the correction of the belief, or the persistence of the image would override the belief. Aristotle draws precisely this inference in the lines that follow: ``either while the fact has not changed and the observer has neither forgotten nor lost belief in the true opinion which he had, that opinion has disappeared, or if he retains it then his opinion is at once true and false. A true opinion, however, becomes false only when the fact alters without being noticed'' (\textit{DA} III.3, 428b4--9). The \textit{phantasma} cannot, therefore, be the bearer of the truth-value; the \textit{doxa} must be. Accordingly, \textit{doxa} is a distinct assertoric act that takes the \textit{phantasma} as its material yet can override or diverge from its content.

\subsubsection*{\textit{The Taking-as-True-or-False} Function: \textit{Doxa}'s Constitutive Operation}

The question that now presses upon the account is how \textit{doxa} arises at all, given that it is neither within our power to summon nor reducible to the \textit{phantasma}'s own testimony. The answer lies in the \textit{taking-as} function---and it is here that the entire architecture of the chain turns. Two distinct operations must be distinguished under this name, and their distinctness is load-bearing for everything that follows. The first is \textit{taking-as}: the presenting of an object under a determinate form together with its affective valence. This is \textit{phantasia}'s accomplishment. The second is \textit{propositional} \textit{taking-as-true-or-false}: the ratification of the \textit{phantasma} as a truth-claim about what is the case. This is \textit{doxa}'s accomplishment. They are related---the propositional ratification operates upon the presentation---but they are structurally distinct, and their distinctness is precisely what the chain requires in order to accommodate the navigation of non-rational animals, the simple appetitive action of rational animals, and the evaluatively complex action through which higher-order emotion enters the sequence.

\textit{Phantasia} is the \textit{kinesis} that arises from the actual exercise of a power of sense (\textit{DA} III.3, 428b11, 429a1--2). The \textit{phantasma} is the settled image that this motion presents to the mind---no more than the formal-epistemic content preserved after the object has withdrawn (the resonant \textit{aisthēma}) together with the affective valence inherited from perception (the \textit{resonant epithymia}). The \textit{phantasma} itself is not a semantic agent; it does not predicate, does not assert. What accomplishes the presenting of this dual content under a determinate aspect---the bear as predator-to-be-avoided, the water as drink-to-be-pursued, the shadow as threat-to-be-fled---is the activity of \textit{phantasia}, through which the \textit{phantasma} comes before the mind as object-with-form-and-valence. The \textit{phantasma} is the medium in which the activity of \textit{phantasia} takes shape.

Two features of \textit{phantasia}'s \textit{taking-as} deserve particular emphasis, for they ground both the possibility of animal action and the structural boundary between \textit{phantasia} and \textit{doxa}. First, the affective valence is not superimposed upon a prior neutral presentation but is constitutive of \textit{phantasia}'s very mode of presenting. \textit{Phantasia} does not first present a neutral image and then tag it with valence in a subsequent cognitive operation; the \textit{resonant epithymia} is co-given with the \textit{resonant aisthēma} as a single disclosure. If this were not so---if affective valence required a separate judgment layered upon neutral perception---non-rational animals could not navigate their world at all, for they possess \textit{phantasia} but lack of \textit{logos} and could not accomplish the layering. Yet Aristotle insists, and the facts of animal life confirm, that the brute does navigate: it has imagination, responds to it, pursues and avoids, acts (\textit{DA} III.11, 434a5--10; III.3, 428a19--24). That it does so is evidence that \textit{taking-as}---including the valence that makes the presentation action-guiding---is pre-propositional, does not require \textit{logos}, and is accomplished by \textit{phantasia} alone.

Second, this \textit{taking-as} is not yet \textit{doxa}. Presenting the bear-\textit{as}-predator-to-be-avoided is not the same function as asserting that \textit{the-bear-IS-a-predator} as a truth-claim about what is the case. The asserting---the propositional \textit{taking-as-true}---is a distinct operation that requires the three conditions Aristotle specifies: belief, conviction, and discourse of reason. What \textit{doxa} adds to \textit{phantasia}'s presentation is neither the \textit{resonant aisthēma} (formal-epistemic content) nor the \textit{resonant epithymia} (the affective valence) which are both supplied by \textit{phantasia}, but the propositional ratification of the \textit{phantasma} \textit{as-true-or-false}. When the brute's \textit{phantasia} presents water-as-drink, appetite responds to the aspectual presentation and the brute drinks; no \textit{doxa} is engaged, and no truth-claim is asserted. When the rational animal's \textit{phantasia} presents water-as-drink, \textit{doxa} may or may not engage; in the simple appetitive case it does not, and the rational animal drinks through the same route as the brute (\textit{MA} 701a32--33). In evaluatively complex cases, \textit{doxa} does engage, and its propositional ratification of an evaluatively loaded presentation articulationally concretizes the generic \textit{resonant epithymia} that \textit{phantasia} has co-presented into the determinate \textit{pathē} that mobilizes higher-order evaluative action.

That \textit{taking-as} and propositional \textit{taking-as-true-or-false} are structurally distinct operations is further demonstrated by the sun passage's phenomenology of divergence. If \textit{phantasia}'s making-present were itself propositional assertion---if presenting the sun as a foot wide were the same function as committing the soul to \textit{the sun is a foot wide}---then the simultaneous holding of a ``one-foot'' \textit{phantasma} with a ``larger than the inhabited earth'' \textit{doxa} would be impossible, for the contrary propositional commitment would cancel the earlier assertion just as one assertion cancels another. But the \textit{phantasma} is not cancelled; it persists (\textit{DA} III.3, 428b2--4). The persistence shows that \textit{phantasia}'s \textit{taking-as} is not in the business of asserting truth at all, and therefore is not suppressed when truth is asserted elsewhere.


\textit{Doxa} arises from this \textit{taking-as}. When \textit{phantasia}, through the \textit{phantasma}, discloses its object as something determinate---as a snake, as a threat, as an enemy approaching in arms---the soul does not merely register the image and move on. The soul takes the disclosure as testimony about what is the case, and this \textit{taking-as-true-or-false} is \textit{doxa}'s originary movement. O'Gorman captures the external grounding of this operation with precision: ``opinion's external grounding---we form opinions about the world about us, and the world about us forms our opinions\ldots. Through \textit{phantasia}, opinion makes judgments about the world of appearances'' (\cite{ogorman_aristotles_phantasia}, 25) that are either true or false. \textit{Phantasia} furnishes the ground; \textit{doxa} ratifies it as true or false. The ratification is neither a separate cognitive achievement that follows the presentation after some interval nor a deliberate weighing of evidence that the soul undertakes upon reflection; it is the soul's involuntary response to the \textit{taking-as} that \textit{phantasia} has already accomplished. \textit{Phantasia} discloses the object as X; the soul, involuntarily, takes X to be the case. This involuntary \textit{taking-as-true-or-false} is what \textit{doxa} articulates. Aristotle marks the asymmetry in a single stroke: ``imagination does not involve opinion based on inference, though opinion involves imagination'' (\textit{DA} III.11, 434a10--11). \textit{Phantasia} can occur without \textit{doxa}---the brute has images, navigates by them, but does not opine---but \textit{doxa} cannot occur without \textit{phantasia}, because the \textit{taking-as} upon which \textit{doxa}'s propositional ratification operates is supplied by \textit{phantasia} alone.

The interface between the two operations is the hinge upon which the entire downstream architecture of the chain depends. \textit{Phantasia} presents, from $A_1$ onward, a \textit{resonant epithymia} constitutive of its making-present---the pleasant-or-painful orientation that Aristotle identifies as co-given with perception (\textit{DA} II.3, 414b4--6). This valence suffices, in simple cases, to move the animal: \textit{resonant epithymia} responds to the aspectually presented pleasant or painful object and appetite fires, as in the drinking case, without any propositional ratification. In evaluatively complex cases, however, the \textit{taking-as} is not itself sufficient for higher-order emotion. When the soul merely imagines something fearful---entertaining the \textit{phantasma} without any present sensible or practical engagement---Aristotle insists that ``we remain unaffected''; it is only when we \textit{take-as-fearful}---when \textit{doxa} propositionally ratifies the aspectual presentation as disclosing a genuine threat---that ``emotion is immediately produced'' (\gk{εὐθὺς πάσχομεν}, \textit{DA} III.3, 427b21--24). The propositional \textit{taking-as-true-or-false}, in other words, is what concretizes the generic \textit{resonant aisthema} and \textit{orexis} that \textit{phantasia} has presented into the determinate \textit{pathē} or emotion which will be analyzed in the next section. \textit{Doxa} is thus not merely one cognitive operation among others but the specific operation through which \textit{phantasia}'s aspectual \textit{taking-as} becomes the propositional \textit{taking-as-true-or-false} that the chain requires in order to generate determinate emotional action.

\subsubsection*{Doxa as Orthogonal Dimension}

The structural consequence is that \textit{doxa} cannot be treated as a fourth cognitive mode parallel to intellection, memory, and discursive \textit{phantasia}. Each mode orients the soul toward the \textit{phantasma}'s content along a distinct temporal and categorial axis. \textit{Doxa}, by contrast, does not orient along any such axis; it ratifies along every axis. One may remember and simultaneously hold a belief about what is remembered; one may deliberate and form opinions about the projected outcomes; one may grasp a universal through intellection and hold a \textit{doxa} about the truth of the proposition apprehended. \textit{Doxa}, therefore, engages upon any orientational mode, adding to each the dimension of involuntary assertoric commitment. It is, accordingly, not a fourth mode but an orthogonal dimension---the committal layer through which any of the three orientations becomes a truth-claim upon which the soul is bound. To characterize \textit{doxa} as an orthogonal committal dimension is to insist that it cuts across the three orientation modes, applying its verdict of truth or falsity to any content they produce. This orthogonality is essential for preserving the autonomy of \textit{phantasia} from \textit{doxa}. The \textit{phantasma} can be generated, manipulated, and combined without any commitment to its truth; \textit{doxa} enters only when the soul takes a stand upon the content that the \textit{phantasma} presents.\footnote{Thomas Rickert's \textit{Ambient Rhetoric} supplies a Heideggerian register in which this orthogonality of \textit{doxa} to \textit{phantasia} finds a contemporary resonance. On Rickert's reading of the Heideggerian disclosure/withdrawal pairing, ``This withdrawing is also a form of supporting; it is the repletion of all that is'' (Rickert 244): every disclosure leaves something unsaid, every showing-forth carries its own holding-back, and the unsaid is not a defect of the disclosure but constitutive of it. The structural homology with the present chapter's thesis is precise at the level of architecture: the \textit{phantasma} presents what \textit{doxa} may or may not ratify, and the unratified content is not thereby annulled --- it persists as available-but-uncommitted, just as the withdrawn dimension persists as constitutive-but-unsaid. Rickert further develops the point through the Heideggerian thing-as-gathering: ``Things `gather' the fourfold ... a thing whose space is carved out by material affordances'' (Rickert 243), where the gathered thing's significance is constituted as much by what it occludes as by what it discloses. The Aristotelian \textit{phantasma}-\textit{doxa} architecture is, on this reading, the cognitive-architectural correlate of disclosure-as-always-also-withdrawal: \textit{doxa}'s committal verdict is a disclosure that simultaneously withdraws the alternative ratifications the same \textit{phantasma} could have received.}

(NOTE TO SYSTEM: see the Rickert footnote attached to the paragraph's final sentence --- we need more research on the disclosure/withdrawal pairing from Rickert and Heidegger. This is very similar to the concept of attunement and is linked to Gibsonian affordances. Every attunement discloses some possibilities for action while occluding others. This also links to Uexküll's tonality: an ``escape tone'' reveals some possibilities for action while concealing others. Rickert's/Heidegger's disclosure/withdrawal, Gibsonian affordance, and Uexküll's tonality all seem to operate on the same principle of disclosing some while concealing others depending on how one is attuned in that moment.)

\subsubsection*{Doxa's Function in the Chain}

Doxa, then, is the soul's moment of commitment: the involuntary propositional ratification, grounded in \textit{phantasia}'s \textit{taking-as} and answerable to how things are, through which appearance becomes truth-claim and generic \textit{resonant epithymia} becomes the determinate mobilization of higher-order emotion. \textit{Phantasia} supplies the aspectual presentation through the \textit{phantasma}; \textit{doxa} binds the soul to the presentation \textit{as-true-or-false}. What \textit{phantasia} discloses aspectually, \textit{doxa} ratifies propositionally; what \textit{doxa} ratifies, emotion mobilizes into determinate conative form; what emotion mobilizes, action completes. The propositional \textit{taking-as-true-or-false} is, in this sense, the hinge of the entire chain for rational animals engaged in evaluatively complex action: the pivot upon which the world's aspectual disclosure becomes the soul's determinate movement. Furthermore, Burke's observation that Aristotle ``looks upon rhetoric as a means that ``proves opposites'' illuminates the orthogonal character of \textit{doxa} (\textit{RoM} 52). The rhetorical capacity to argue for and against the same proposition presupposes that the propositions themselves---grounded in \textit{phantasmata}---are  neutral with respect to truth-value until \textit{doxa} commits the soul to one side. Heidegger notes this poignantly when he remarks that ``ultimately, rhetoric has its sight on \gk{κρίσις} (\textit{krisis} / ``judgment''), view-construction, a definite decision in the sense of a \gk{δόξα} (\textit{doxa})'' (110). Here \textit{krisis}---the formation of a definite view---is presented as the \textit{telic} function of rhetoric, and \textit{doxa} is what \textit{krisis} aims at: for, as Heidegger clearly observes, ``[v]iews are cultivated, renewed, established, hardened in speaking'' and rhetoric, as such, ``is nothing other than the \textit{interpretation of being-there with regard to the basic possibility of speaking-with-one-another}'' (95). The judgment is precisely the act of \textit{taking-as-true}, the involuntary ratification through which appearance becomes truth-claim. Thus, the orthogonality of \textit{doxa} is not merely a structural feature of individual psychology but a condition of the possibility of rhetoric itself: only because \textit{phantasia} operates independently of \textit{doxastic} commitment can the rhetor construct opposing \textit{phantasmata} and invite the audience to form a definite view of one \textit{doxastic} stance over another.\footnote{Burke's identification of rhetoric as ``the art of persuasion, or a study of the means of persuasion available for any given situation'' (Burke 47) operates on precisely this orthogonality: persuasion is possible because the same \textit{phantasma} can be ratified differently by different \textit{doxai}, and the rhetorical work consists in cultivating the dispositional ground from which one ratification rather than another becomes likely. Burke's notion of identification---``in \textit{acting together}, men have common sensations, concepts, images, ideas, attitudes that make them consubstantial'' (RoM 21)---specifies the social-ontological mechanism by which collective \textit{doxa}-formation occurs. Gross extends Burke's identification doctrine into a contemporary rhetorical-affordance theory where ``rhetoric, as a humanistic form of affordance theory, posits a situation that is not neutral but is rather `persuasive,' threatening and promising with respect to our human being-in-the-world'' (Gross 2017, pp.~22--23). The orthogonality of \textit{doxa} is, on this combined reading, the cognitive-architectural condition for both the Aristotelian art of rhetoric and its contemporary affordance-theoretic articulation.}

\subsubsection*{Conclusion: The Architecture of A\textsubscript{3} and Its Transition into Action}

The completed cognitive actuality at $A_3$ is not a single homogeneous state but a multi-dimensional achievement in which several distinct operations converge upon the \textit{phantasma} stabilized at $A_2$. The determinate, available, representational image that \textit{phantasia} has stabilized from the \textit{resonant kinesis} of perception now stands ready for the full range of cognitive and \textit{orectic} operations that characterize the life of the rational soul. The rational cognitive apparatus engages the \textit{phantasma} along three orientational axes: the atemporal axis of intellection, which addresses the \textit{phantasma} as bearing a universal; the retrospective axis of memory, which addresses the \textit{phantasma} as a likeness of what has elapsed; and the discursive axis of deliberative \textit{phantasia}, which addresses multiple \textit{phantasmata} as composable into either projected future possibilities (the practical sub-domain) or contemplated unchangeable truths (the speculative sub-domain). These three modes differ from one another in intentional direction and temporal orientation, but they share a common structural role as non-committal engagements through which the \textit{phantasma}'s content is taken up and processed without yet being ratified as true. Alongside the three, and structurally distinct from them, operates the committal dimension of \textit{doxa}, which engages orthogonally upon any orientational mode and adds to each the propositional \textit{taking-as-true-or-false} through which the soul is bound involuntarily to a truth-claim about what that mode has disclosed.

This four-dimensional architecture is organized, throughout, by the two-level structure of taking-as that governs $M_2 \rightarrow A_3$ as a whole. \textit{Phantasia}'s \textit{taking-as}---the presenting of the object with a \textit{resonant aisthēma} aspect together with its \textit{resonant epithymia}---makes the \textit{phantasma} cognitively consequential and supplies the material upon which every operation here proceeds. It is pre-propositional, does not require \textit{logos}, and is accomplished by \textit{phantasia} alone. \textit{Doxa}'s propositional \textit{taking-as-true} is \textit{logos}-dependent and operates exclusively in rational animals, ratifying what \textit{phantasia} has presented as a truth-claim about what is the case. The three orientational modes operate at both levels: each engages the \textit{phantasma} in its distinctive cognitive direction, and each can be ratified \textit{doxastically} when the rational animal commits the soul to the truth of what the mode discloses. 


The structural consequence for the larger chain follows directly. $A_3$ serves, conditionally, as the unmoved originator of the subsequent motion at $M_3 \rightarrow A_4$, and the condition is that $A_3$'s content disclose a realizable good. Where the completed actuality is the contemplation of an unchangeable truth---the speculative case whose object is and cannot be otherwise---no realizable good is disclosed and the chain terminates at $A_3$; speculative thought, as Aristotle observes, moves nothing (\textit{DA} III.10, 433a13--15). Where the completed actuality bears on a practical possibility---a remembered pattern informing present deliberation, an anticipated future evaluated against a single standard, an intellected universal entering a practical syllogism, a \textit{doxa} ratifying an evaluatively complex aspect---the chain continues, and $A_3$ functions as the apparent good that moves \textit{orexis}. The passage from $A_3 \to A_4$ represents the moment at which the soul moves from having an image to doing something with that image---from passive presentation to active engagement. Aristotle's insistence that what originates movement is both pre-eminently and primarily soul, and that what is not itself moved cannot originate movement in another (\textit{De Anima} III.10, 433a13--18: ``both thought and appetite are capable of originating local movement\ldots\ the object of appetite is the stimulant of practical thought''), establishes the \textit{orectic} dimension as the driving force of the $A_3 \to A_4$ transition. The \textit{phantasma} does not merely present a cognitive content; it engages the soul's desire, and it is desire---\textit{orexis}---that moves the animal to action. Thus $A_3$, understood as the fully constituted \textit{phantasma} together with its \textit{doxastic} overlay, is the proximate ground of practical motion: the soul acts because it desires, and it desires because the \textit{phantasma} presents an object of desire under a \textit{doxastic} commitment to its goodness or badness. The object of practical desire, Aristotle stipulates, is ``the real or apparent good'' (\textit{to prakton agathon \=e phainomenon}, \textit{De Anima} III.10, 433a27--29); this ``apparent good'' component is what allows the \textit{phantasma} at $A_3$ to function as the ratifying-articulation that moves \textit{orexis} even when the \textit{prakton} is only apparently good. Aristotle identifies nature itself as a principle of being-moved and of being-at-rest in that to which it belongs primarily and in virtue of itself (\textit{Physics} II.1, 192b20--23: ``nature is a principle or cause of being moved and of being at rest in that to which it belongs primarily, in virtue of itself and not accidentally''); similarly, the \textit{phantasma} is the principle of cognitive and practical motion in the rational soul, that from which all higher activity originates and to which it returns. The passage from $A_3 \to A_4$ accordingly marks the threshold between the constitutive and the operative dimensions of the actualization chain. $A_3$ is constitutive: it brings the \textit{phantasma} into being as a determinate, available, representational image. 

$A_3$ accordingly occupies the pivotal location in the chain. Everything prior---the perceptual motion from $A_0$ through $A_1$, the phantastic motion from $A_1$ through $A_2$---converges upon it, and everything subsequent---the \textit{orektikon} motion at $M_3 \rightarrow A_4$, the completed action at $A_4$---depends upon what it produces. It is the point at which the world's disclosure, carried by the \textit{phantasma} through the residual motion of \textit{phantasia}, is taken up cognitively by the rational apparatus and, in rational animals engaged with evaluatively complex content, ratified propositionally through \textit{doxa} as the ground upon which determinate emotional action will unfold. What the soul has perceived and retained becomes, at $A_3$, what the soul takes to be the case; and what the soul takes to be the case, where it discloses a realizable good, becomes what the soul is moved to enact.

Therefore, we can conclude this subsection with a refined formulation: \textit{phantasia}, in its \textit{taking-as-something} function, is the bodily-disclosive ``making-present'' through which the \textit{phantasma} presents the \textit{eidos} of the perceived particular with is \textit{resonant epithymia}. It is pre-propositional in not yet involving  \textit{logos}-articulation, but it is structurally \textit{eidos}-disclosing---the \textit{phantasma} presents the being \textit{as} such-and-such precisely because it bears the \textit{eidos}. \textit{Taking-as-something-true-or-false} is the \textit{logos}-dependent \textit{phasis} (``yes-saying'') through which \textit{doxa} ratifies the disclosure as a truth-claim. 

\footnote{Heidegger reads \textit{doxa} through the lens of his analysis of discourse and view-construction, observing that rhetoric has its sight on \textit{krisis} --- a definite decision in the sense of a \textit{doxa} --- and that this decision is cultivated through discourse itself (\textit{BCAP}, pp.~110, 124). Heidegger's analysis at \textit{BCAP} p.~124 frames Aristotle's ethical analysis as ``a sharp contrast between \textit{legein} about ethics-related problems and real philosophizing,'' opposing the sophists' and moralizers' merely-public orientation by re-locating ethical practice in the agent's settled comportment: the triad of conditions on the one acting --- \textit{eidos} (knowing, oriented toward \textit{kairos} with \textit{phron\=esis}); \textit{proairoumenos} (acting from out of himself); \textit{bebaios kai ametakin\=etos echōn} (``stable and not to be brought out of composure'') --- specifies what it means for the rational animal to hold a \textit{doxa} positively, rather than merely as one opinion among many in public exchange. Accordingly, \textit{doxa} is not merely a private psychological event but a publicly oriented stance: even when the soul takes-as-true in the privacy of inner deliberation, it does so in a manner that is structurally homologous with the public expression of opinion in rhetorical contexts.}

\inlinenote{Papachristou's insitence on distinguishing the grades of phantasia---sensory, deliberative, adn potentially a third kind associated with rational contemplation---supports this orthogonality, for each grade produces content that is subject to doxastic evaluation yet is not itself a form of belief (9-10).}


\end{document}